\documentclass[12pt]{article}
\usepackage{amsthm,amssymb,amsmath}
\usepackage{bm,bbm}
% \usepackage{mathrsfs}
\usepackage{mathtools,etoolbox,xcolor}
% \usepackage{mathtools}
\usepackage{enumitem}
\usepackage[margin=1in]{geometry}
\usepackage{authblk}
\usepackage{natbib}
% \usepackage[style=authoryear, natbib=true, backend=biber]{biblatex}
\usepackage{booktabs,siunitx}
% \usepackage{multirow}
\usepackage{caption,subcaption}
% \usepackage{xr}
\usepackage{tikz}
\newtheorem{theorem}{Theorem}
\newtheorem{definition}[theorem]{Definition}
\newtheorem{corollary}[theorem]{Corollary}
\newtheorem{proposition}[theorem]{Proposition}
\newtheorem{lemma}[theorem]{Lemma}
\newtheorem{conjecture}{Conjecture}
\renewcommand{\P}{\mathbbm{P}}
\newcommand{\E}{\mathbbm{E}}
\newcommand{\var}{\text{Var}}
\newcommand{\corr}{\text{corr}}
\newcommand{\munull}{\mu_{0}}
\newcommand{\tausqnull}{\tau^2_{0}}
\newcommand{\taunull}{\tau_0}
\newcommand{\tausqhat}{\hat\tau^2}
\newcommand{\mumle}{\hat\mu_{MLE}}
\newcommand{\tausqmle}{\hat\tau^2_{MLE}}
\newcommand{\tausqmlered}{\hat\tau^2_{MLE,R}}
\newcommand{\mustar}{\mu_\ast}
\newcommand{\tausqstar}{\tau^2_\ast}
\newcommand{\comment}[2][red]{%
 \iftoggle{commenttoggle}{%
  {\normalsize{\color{#1}{#2}}\normalsize}%
 }{%
 }%
}
\graphicspath{{./figs/}} 
\newtoggle{commenttoggle}
\togglefalse{commenttoggle}
\mathtoolsset{showonlyrefs,showmanualtags}

\title{Bartlett adjustment for random effects meta-analysis of event outcomes with small samples}
% \date{}
\author[1]{Haben Michael}
\affil[1]{University of Massachusetts, Amherst, MA (hmichael@math.umass.edu)}
\date{}                   
\begin{document}
\maketitle

Abstract. ...
\textsc{Keywords: Meta-analysis, Bartlett correction, Generalized linear mixed models}


1. introduction: (skip generic metaanalysis material). since stijnen
2010, glmm have become the preferred model for random effects
meta-analysis based on count data. cochran etc have recommended
... . cite elsewhere for advantages over inverse variance model. while
it avoids the normal approximation, inference is still based on a chi
squared asymptotic expansion. therefore as before in the IV case the
many smaller meta-analyses that are carried out are suspect. bartlett
correction is a way to improve the efficiency of inference. when
assumptions are met, allows rate of convergence to jump from 1/n to
1/n2. the key assumption in this setting is the gaussianity of the
random effect.

relevant literature. 1. bartlett corrections applied to glms. what
about glmms like this? 2. hanada paper. main difficulty is
computational. with careful design we argue that these difficulties
can be overcome and consumer hardware is sufficient for practical use.

paper outline. in methods, give stifjnen's glmm approach as
background, then describe algo. next look at sims, focusing on
robustness. for data analysis walk through the 4 stijnen analyses.


2. background.

2a. notation for events, trials. 4 types of analyses: single arm and
two arm trials, proportions and incidence/rates.

2a-1 single arm, proportion: binomial-prop
2a-2 two arm, proportion: hyperg
2a-3 single arm, incidence: poisson
2a-4 two arm, incidence rate ratio: binomial-irr

in all cases the likelihood has the form: int condpmf * phi. condpmf = , parameters= (used loosely--need to distinguish from mu/tau2).

give LR test.

2b. bartlett correction

bartlett correction. write out expansion of likelihood. odd powers of root n vanish. identify c as the 1/n term. so dividing by 1+c then removes the mean bias to order 1/n,
\begin{align}
E(...) =.
\end{align}
the obesrvation of ((lawley? or bartlett)) is that this mean bias
reduction in fact removes the 1/n bias from ((ref LR expansion)) as
well, i.e., the chi squared approximation. in the normal course with a well-specified model the resulting asymptotic rate is 1/n2. may be used with profile curve to obtain CIs and p-vals. ((cite to sim/data analysis.))

background: lawley formulas. k terms.

an important exception to the ``normal course''  is lattice data such as considered here. expansion ((lr expansion)) may not be justified. literature unclear. bartlett corrections for discrete glm's have been developed.



2c. proposed method.

to compute the lawley formulas for ((ref likelihood)) we approximate the integral by a sum. since all the likelihoods involve integrating against a gaussian, gauss-hermite approximation is a natural choice, replacing the integral ((ref)) with a weighted sum
the weights are ... . ((how much detail to give here, maybe just leave it as a sum approximation.)) let r denot eth emultivector, the derivative of the log likelihood wrt to r is
\begin{align}
\end{align}
and its expectation
\begin{align}
\end{align}
may be approximated by another quadrature
\begin{align}
\end{align}
expression ((...)) approximates k terms of the form ((inner only)). k terms of the form ((inner and outer)) may be approximated by 
\begin{align}
\end{align}
where as before s is a multivector of indices.

2d. further implementation details.

chain rule: differentiate wrt pmf sums. get 35 k terms in terms of pmfsums (give an example or two).

differentiate pmfs using recurrences in terms usual pmfs to allow fast, stable routines to be used.

some specific issues: support for poisson. speed issues with hyperg.


3. simulation

goals: robustness/gaussianity check. discreteness issue.

4. data analysis.

4a. proportions

4b. incidence


\bibliographystyle{chicago}
\bibliography{../ma-equivalence/ma-equivalence.bib}
% \printbibliography

\end{document}


